% =========================================================
% APPENDIX A — Affine–Projection Constants
% =========================================================

\section{Affine–Projection Constants}
\label{app:A}

This appendix records explicit constants for the blockwise
short–interval variance bounds used in
Theorem~\ref{thm:variance} (variance regime, Part~1) and in
Step~\ref{step:residual-variance} of Part~2.

\subsection{Set-up and Frequency Partition}

Write the smoothed remainder along $\Re s=\sigma$ as a Fourier–type
superposition over zeros:
\[
H_\sigma(t)=\sum_{\rho}\frac{e^{(\beta-\sigma)t}}{|\rho|^{\alpha}}\cos(\gamma t),
\qquad
\eta_\rho:=\beta-\sigma .
\]
Differentiating gives
\[
H'_\sigma(t)=\sum_{\rho}\frac{e^{\eta_\rho t}}{|\rho|^{\alpha}}
\bigl(\eta_\rho\cos(\gamma t)-\gamma\sin(\gamma t)\bigr).
\]
Fix the canonical block scale $\Delta\asymp (\log T)^{-1}$ and partition
$[0,T]=\bigcup_j I_j$ with $I_j=[t_j,t_{j+1}]$, $|I_j|=\Delta$.
As in Remark~\ref{rmk:truncation}, at the explicit-formula value $\alpha=1$ the
derivative field is taken in truncated form, with frequencies
$|\gamma|\le X:=\Delta^{-1}\asymp\log T$.  In complex notation
\[
H_\sigma'(t)=\sum_{|\gamma|\le X} b_\rho\, e^{(\eta_\rho+i\gamma)t},
\qquad |b_\rho|\asymp |\rho|^{-\alpha}|\gamma|,\quad \eta_\rho=\beta-\sigma .
\]
The argument below uses only the \emph{global} second moment of this truncated
field; its leading term is unconditional, and the sole arithmetic input is the
mild ordinate-spacing Hypothesis~\ref{hyp:spacing}~(S), far weaker than any
large-sieve minimal-gap requirement.

\subsection{Blockwise Affine Projection and a Mean–Variance Lemma}

Let $P_j$ denote the $L^2(I_j)$ projection onto affine functions
$\{a+bt\}$ and set $R_j=(\mathrm{Id}-P_j)H_\sigma$.  Then $R_j'=H_\sigma'-b_j$,
where $b_j$ is the (constant) slope of the affine fit.  We have the blockwise
stability
\begin{equation}
\int_{I_j} |R'_j(t)|^2\,dt
\;\le\;
c_{\mathrm{aff}}^{-1}\int_{I_j} |H'_\sigma(t)|^2\,dt,
\label{eq:aff-proj}
\end{equation}
for an absolute constant $c_{\mathrm{aff}}\in(0,1)$.  Indeed, the regression slope
$b_j=\frac{12}{\Delta^3}\int_{I_j}(t-\bar t)\,H_\sigma\,dt
=\frac{12}{\Delta^3}\int_{I_j}\kappa_j(t)\,H_\sigma'(t)\,dt$ for a kernel
$\kappa_j$ with $|\kappa_j|\le\Delta^2$, so by Cauchy--Schwarz
$\Delta\,b_j^2\le C\int_{I_j}|H_\sigma'|^2$; combined with
$\int_{I_j}|H_\sigma'-b_j|^2=\int_{I_j}|H_\sigma'|^2-\Delta\,b_j(2m_j-b_j)$,
where $m_j=\Delta^{-1}\!\int_{I_j}H_\sigma'$ is the block mean of $H_\sigma'$ and
$\Delta\,m_j^2\le\int_{I_j}|H_\sigma'|^2$, this yields \eqref{eq:aff-proj} with an
absolute constant.  (We make no use of any specific value of $c_{\mathrm{aff}}$.)

\subsection{Reduction to a global second moment}

Because the blocks tile $[0,T]$, \eqref{eq:aff-proj} gives
\begin{equation}
\sum_j\int_{I_j}|R_j'|^2
\;\le\; c_{\mathrm{aff}}^{-1}\sum_j\int_{I_j}|H_\sigma'|^2
\;=\; c_{\mathrm{aff}}^{-1}\int_0^T|H_\sigma'(t)|^2\,dt,
\label{eq:reduce-global}
\end{equation}
so it suffices to bound the global second moment of the truncated field.

\subsection{Diagonal main term}

Expanding the square,
\[
\int_0^T|H_\sigma'|^2\,dt
=\sum_{|\gamma|\le X}|b_\rho|^2\!\int_0^T\! e^{2\eta_\rho t}\,dt
\;+\;\sum_{\rho\ne\rho'} b_\rho\bar b_{\rho'}\!\int_0^T\! e^{(\eta_\rho+\eta_{\rho'})t}e^{i(\gamma-\gamma')t}\,dt .
\]
In the variance regime $\eta_\rho=\beta-\sigma\le0$.  An on-line zero
($\eta_\rho=0$) contributes $\int_0^T e^{2\eta_\rho t}dt=T$ to the diagonal;
an off-line-left zero ($\eta_\rho<0$) contributes $\le 1/(2|\eta_\rho|)=O(1)$.
Hence the diagonal equals
\[
T\!\!\sum_{\substack{0<\gamma\le X\\ \beta=\sigma}}\!\!|b_\rho|^2\,(1+o(1))
=T\!\!\sum_{0<\gamma\le X}\!\frac{\gamma^2}{|\rho|^{2\alpha}}\,(1+o(1)),
\]
the off-line-left zeros giving a lower-order contribution.  By
Riemann--von~Mangoldt, $dN(V)=\tfrac{1}{2\pi}\log\tfrac{V}{2\pi}\,dV+O(dV)$, so
with $|\rho|\asymp\gamma$,
\[
\sum_{0<\gamma\le X}\frac{\gamma^2}{|\rho|^{2\alpha}}
\asymp\int_1^X V^{2-2\alpha}\,\frac{\log V}{2\pi}\,dV
=\begin{cases}\dfrac{X^{3-2\alpha}\log X}{2\pi(3-2\alpha)}(1+o(1)), & \alpha<\tfrac32,\\[6pt] O(1), & \alpha>\tfrac32.\end{cases}
\]
At the explicit-formula value $\alpha=1$ this is
$\frac{X\log X}{2\pi}(1+o(1))=\frac{\log T\log\log T}{2\pi}(1+o(1))$ since
$X=\Delta^{-1}\asymp\log T$, so the diagonal is
\begin{equation}
T\sum_{0<\gamma\le X}\frac{\gamma^2}{|\rho|^{2\alpha}}
=\frac{1}{2\pi}\,T\log T\log\log T\,(1+o(1)).
\label{eq:diag-main}
\end{equation}
This is the source of both the $\log\log T$ factor and the criticality of
$\alpha=1$ (for $\alpha>1$ the sum converges and the bound improves to $O(T)$).

\subsection{Off-diagonal: negligibility via zero separation}

Using $|e^{(\eta_\rho+\eta_{\rho'})t}|\le1$ and
$\bigl|\int_0^T e^{i(\gamma-\gamma')t}dt\bigr|\le 2/|\gamma-\gamma'|$, the
off-diagonal is bounded by
\[
\Bigl|\sum_{\rho\ne\rho'}\cdots\Bigr|
\le 2\!\!\sum_{\substack{|\gamma|,|\gamma'|\le X\\ \gamma\ne\gamma'}}\!\!
\frac{|b_\rho||b_{\rho'}|}{|\gamma-\gamma'|}
\le\frac{2}{\delta_{\min}(X)}\Bigl(\sum_{|\gamma|\le X}|b_\rho|\Bigr)^2 ,
\]
where $\delta_{\min}(X)$ is the least gap between distinct ordinates with
$|\gamma|\le X$.  The diagonal carries the full factor $T$ from $\int_0^T$, so
the off-diagonal is negligible as soon as $\delta_{\min}(X)^{-1}$ is small
relative to $T$.  This is guaranteed by the following mild input, which we state
explicitly as it is the \emph{only} arithmetic hypothesis entering the variance
bound.

\begin{hypothesis}[Mild ordinate spacing, \textup{(S)}]\label{hyp:spacing}
There is an absolute $c_0>0$ such that distinct ordinates $\gamma\ne\gamma'$ of
$\zeta$ with $|\gamma|,|\gamma'|\le V$ satisfy $|\gamma-\gamma'|\gg V^{-c_0}$;
equivalently $\delta_{\min}(V)\gg V^{-c_0}$.
\end{hypothesis}

\noindent
Hypothesis (S) is far weaker than the Riemann Hypothesis or the pair-correlation
conjecture: it only forbids distinct ordinates from approaching faster than any
fixed polynomial rate, and is consistent with all known zero computations
(the closest recorded pairs, e.g.\ Lehmer pairs, are separated by $\gg 10^{-5}$,
vastly larger than the required $V^{-c_0}$ at the relevant heights $V\asymp\log T$).
What \emph{can} be proved toward (S), and the precise obstruction to a complete
proof, are recorded in \S\ref{subsec:toward-S} below; (S) is no easier than the
known unconditional small-gap problem for $\zeta$-ordinates, and we isolate it as
the single open input.  Under (S),
$\delta_{\min}(X)^{-1}\ll X^{c_0}=(\log T)^{c_0}$ and, since
$\sum_{|\gamma|\le X}|b_\rho|\ll N(X)\ll X^{1+o(1)}$, the off-diagonal is
$\ll(\log T)^{c_0+2+o(1)}$ --- \emph{polylogarithmic} in $T$ --- hence negligible
against the diagonal \eqref{eq:diag-main}.  The point is structural: the
truncation height $X\asymp\log T$ is so small relative to the integration length
$T$ that, granting (S), the cross terms cannot compete; no minimal-gap or
large-sieve estimate for the (clustered) ordinates is invoked.

\subsection{Toward Hypothesis (S)}\label{subsec:toward-S}

We record what can be proved toward (S) and isolate the obstruction.

First, in the variance regime at $\sigma=\tfrac12$ the hypothesis simplifies.
If every zero satisfies $\beta\le\tfrac12$, then by the functional equation
together with the reflection $\zeta(\bar s)=\overline{\zeta(s)}$, each zero
$\rho=\beta+i\gamma$ is paired with $1-\bar\rho=(1-\beta)+i\gamma$, also a zero, of
real part $1-\beta\ge\tfrac12$; both can satisfy $\beta\le\tfrac12$ only if
$\beta=\tfrac12$.  Hence \emph{at $\sigma=\tfrac12$ the variance regime forces
all zeros onto the critical line}, and distinct ordinates correspond to distinct
zeros with $|\rho-\rho'|=|\gamma-\gamma'|$.  For such zeros we have a conditional
separation.

\begin{proposition}[Conditional polynomial separation]\label{prop:cond-sep}
Let $\rho_1=\tfrac12+i\gamma_1,\ \rho_2=\tfrac12+i\gamma_2$ be distinct on-line
zeros with $|\gamma_j|\le V$, set $d=|\gamma_1-\gamma_2|<1$ and
$s_0=\tfrac12+i\tfrac{\gamma_1+\gamma_2}{2}$, and let $A$ be an absolute exponent
with $|\zeta(s)|\ll V^{A}$ on $\{|s-s_0|\le1\}$ \textup{(}the functional equation
and convexity give polynomial growth in the height on any fixed-width strip, so
such an absolute $A$ exists\textup{)}.
If there is a point $c$ with $|c-s_0|\le\tfrac12 V^{-(A+B)/2}$ and
$|\zeta(c)|\ge V^{-B}$, then $d\ge V^{-(A+B)/2}$.
\end{proposition}

\begin{proof}
Jensen's formula for $\zeta$ on the disc $\{|s-c|\le1\}$ gives
\[
\sum_{|z_k-c|\le1}\log\frac{1}{|z_k-c|}
=\frac{1}{2\pi}\!\int_0^{2\pi}\!\log\bigl|\zeta(c+e^{i\theta})\bigr|\,d\theta
-\log|\zeta(c)|
\;\le\; A\log V+B\log V,
\]
the sum over zeros $z_k$ of $\zeta$ in the disc.  Every term is nonnegative, and
$\rho_1,\rho_2$ alone contribute at least $2\log\frac{1}{d/2+|c-s_0|}$ (as
$|\rho_j-c|\le d/2+|c-s_0|$).  Hence $d/2+|c-s_0|\ge V^{-(A+B)/2}$, and with
$|c-s_0|\le\tfrac12 V^{-(A+B)/2}$ this yields $d\ge V^{-(A+B)/2}$.
\end{proof}

\begin{remark}[The obstruction and the status of (S)]\label{rmk:obstruction}
Proposition~\ref{prop:cond-sep} reduces (S) to producing a point $c$ within
$V^{-c_0}$ of the cluster at which $|\zeta|\ge V^{-B}$.  This is exactly what is
\emph{not} available unconditionally: near a cluster of the $O(\log V)$ zeros in
a unit disc, Cartan's lemma bounds the small-value set of $\zeta$ only up to an
exceptional radius of size $O(1)$, not $V^{-c_0}$, so one cannot in general find a
non-small point as close as $V^{-c_0}$ to the cluster.  In this sense (S) is no
easier than the unconditional small-gap problem for $\zeta$-ordinates --- a
recognized hard problem.  Lehmer's phenomenon shows the gaps can be very small,
while random-matrix heuristics predict the smallest gap up to height $V$ to be
only polynomially small, comfortably within (S) for a suitable $c_0$.
We stress that (S) is more than what the variance bound needs: (S) makes the
off-diagonal \emph{polylogarithmic}, but for it merely to be \emph{negligible}
against the diagonal $\asymp T\log T\log\log T$ it suffices that
$\delta_{\min}(\log T)\gg T^{-1}(\log T)^{O(1)}$ --- i.e.\ that no two distinct
ordinates below height $\log T$ lie within $T^{-1}(\log T)^{O(1)}$ of each other,
a vastly weaker requirement, overwhelmingly satisfied since those ordinates are
low-height zeros with gaps $\asymp 1/\log\log T$.  We therefore leave (S) as the
single, mild, numerically unchallenged arithmetic input.
\end{remark}

\subsection{Conclusion}

Combining \eqref{eq:reduce-global}, \eqref{eq:diag-main} and the off-diagonal
bound (the latter under Hypothesis~\ref{hyp:spacing}) gives the global variance
bound
\begin{equation}
\sum_j\int_{I_j}\!|H'_\sigma(t)|^2\,dt
\;\le\; C_0(\sigma,\alpha)\, T\log T\log\log T,
\qquad C_0=\tfrac{1}{2\pi}+o(1)\ \ (\alpha=1),
\label{eq:globalC0}
\end{equation}
and, with \eqref{eq:aff-proj}, the form used in
Step~\ref{step:residual-variance} of Part~2:
\begin{equation}
\sum_j\int_{I_j}\!|R'_j(t)|^2\,dt
\;\le\; C_0'(\sigma,\alpha)\, T\log T\log\log T,
\qquad C_0' = c_{\mathrm{aff}}^{-1} C_0.
\label{eq:residC0prime}
\end{equation}

\begin{remark}[About lower bounds]
Only the \emph{upper} variance bound \eqref{eq:globalC0} is needed for
Theorem~\ref{thm:variance} and for the residual control in Part~2.
Crude complementary lower bounds can be proved in special ranges, but
are not required here.
\end{remark}

\subsection{Numerical Verification}
The constants $c_{\mathrm{aff}}$, $C_0$, and $C_0'$ have been checked in the
supplementary notebooks (Part~3) by blockwise evaluation on synthetic signals
and Odlyzko windows; see the repository cited in Section~\ref{sec:numerics}.
